
\renewcommand{\thechapter}{5}

\chapter{Conclusions and Future Work}\label{chap:conclusion}

The work I present in this thesis aims to enable reliable, scalable, and reproducible
experimentation and research for the next generation of random testing tools. We were
able to undertake experiments that would otherwise have taken weeks to setup in a
matter of hours and days due to the combination of tools and techniques I've worked
during the development of this thesis.

We are in the midst of a paradigm shift in the development of software. Program
synthesis from natural language specifications via generative AI has given rise to a
proliferation of code generation. The leverage one gains from such generation, however,
is bounded by their ability to verify it, putting verification of correctness to the
center of this debate. I wrote this thesis in the midst of this shift, and I have
witnessed the rising interest in random testing during these last 5 years.

Property-Based Testing is particularly interesting within this new paradigm of code
synthesis because it emphasizes specifications as a central concept. Given a complete
specification of a system, the futurist in me wants to claim we could synthesize the
system from scratch. In addition to developing property-based testing tools themselves,
I would like to use random testing for autonomous systems optimization, translation,
and synthesis; develop tools and theories advancing these layers on top of testing. I
believe, however, we are far from this reality, not due to model capabilities, but
rather due to our lack of understanding what makes property-based testing effective.

During my thesis work, I have worked on building the infrastructure for collecting
metrics and measuring statistics, devising which metrics to collect and statistics to
compute and empirically validating them. Today, we can measure the quality of a
generator produced by an algorithm for an ETNA workload; can we go further to
understand the quality of a generator written for an arbitrary project using PBT? Can
we automatically tune internal parameters of a generator as needed? We produced
measures for evaluating shrinking strategies, can we generalize them, use them to find
better algorithms across different ecosystems? Our measurements treat mutation-based
random testing as a direct alternative to blackbox generation-based random testing,
what kind of fine granular metrics can we devise for better understanding the
effectiveness of the former? Finally, how can we leverage Programmable PBT to augment
the current property runners, can we have fault localization in PBT, can we switch
between symbolic and concrete inputs, can we parallelize work only when needed?

We may be able to do all of these things, devise metrics, develop new techniques and
algorithms only if we are able to evaluate them properly. I plan to horizontally grow
ETNA to as many languages, ecosystems and workloads as possible; to make
experimentation across the board as seamless as possible, evaluate any and all features
of a PBT framework, supercharge the development for the future of property-based
testing.
